\section{Evaluation design and inference rules}
\label{sec:v2evaluation}
The evaluation separates implementation acceptance from scientific conclusions.
Proof failures, inconsistent traces, broken semantic controls, isolation bypasses,
and incorrectly generated evidence fail acceptance. Unfavorable or underpowered
model outcomes would narrow the empirical claim rather than fail a gate that
requires a favorable result. Missing model outcomes are reported as missing,
not replaced with scripted proposals carrying checkpoint metadata.

\subsection{Versioned interfaces and retained evidence}
All v1 commands and fixtures retain their meaning. The v1 reference evidence is
bound to commit \texttt{12b5579}; its source hashes are checked against that
revision, not silently rewritten to the current tree. A fresh v1 experiment run
also exercises backward compatibility after the v2 changes. The supplied
\texttt{tips/*.tex} hashes remain unchanged. V2 evidence lives separately and
records a versioned protocol manifest, source and binary hashes, proof/toolchain
identifiers, host and kernel information, and the predeclared model settings.

The \texttt{run-v2} command has smoke and publication profiles. Smoke uses one
repetition per actor population; publication uses three. Both run the same fault
cases, proof checks, semantic controls, and mutation language. The profiles are
not different safety policies. \texttt{analyze-v2} validates the resulting traces
and evidence before optional model-campaign analysis. The CPU gate uses smoke,
while the recorded host report uses publication repetitions. Publication here
names a workload profile, not a claim that a model study was completed.

Invocations and responses are recorded with role, request, monotonic timestamps,
and observed effects read independently from the sink store. Fault histories
have at most one outstanding RPC per service and are checked by a separately
implemented Python specification. The checker reconstructs approvals, issuance,
reservations, issuer withdrawal, sink acceptance, effect application, and
settlement. It compares predicted sink effects to the independent effect log,
including contents and destinations rather than favorable summary flags.

Concurrent races are treated differently. Arbitrarily ordered RPC responses do
not constitute a valid sequential linearization. The race harness starts
independent broker writers against the same store, releases them with a barrier,
and checks final conservation and exact effects. The delayed-fence experiment
also has overlapping calls and uses transaction-lock observation plus independent
post-acknowledgment effects. These tests do not claim an exhaustive concurrent
history checker. On a sequential-history failure the harness retains the first
failing prefix; that is prefix minimization, not a globally minimal trace.

\subsection{CPU faults and adversarial controls}
The fault suite covers useful transfers, releases, and edits; duplicate and
altered requests; actor substitution; a missing required issuer; expired
issuance; admission and fenced withdrawal; independent issuer semantics;
renewal with stale approvals; stale document revisions; lost replies; process
death on both sides of commit; sink restart; uncertain settlement; and first
delivery after fencing. Crash hooks terminate the actual service process before
or after its durable commit. Recovery starts another process against the same
store and retries the identity.

Two controls are intentionally harmful. The unsupported non-idempotent service
duplicates a retried delivery. The specification-omission case commits an
authorized but incorrect document. Acceptance requires observing these failures,
because a checker that cannot detect them would provide weak ground truth.
They are labeled by their cause: a violated service contract versus a task
property absent from permission. Neither is suppressed to obtain a zero-failure
headline.

Ten enforcement mechanisms are individually omitted in an abstract protocol
model: cumulative budget, payload binding, issuer conjunction, epoch freshness,
expiry, revision checking, deduplication, durable accounting, uncertainty
quarantine, and sink fencing. All 45 pairs are also searched. The search space is
every subsequence of ten declared regression templates, sorted by length. State
updates and the harm oracle are distinct in the model. Each witness is shortest
within that finite language. It is not necessarily shortest among arbitrary
protocol executions, nor evidence that all possible interacting failures have
been found.

The abstract mutations are not alternate versions of the verified Rust core.
They test the roles of mechanisms in the protocol argument. Likewise, a pair
whose witness already exploits one omitted mechanism does not establish a new
synergistic interaction. The evidence retains every omission set, witness, harm
type, template, and searched-trace count so readers can identify that limitation.
The complete abstract model also admits a useful task and has no harm within
the searched language.

\subsection{Controllers and availability measurements}
The declared comparison has five controllers: no enforcement, stateless per-call
rules, learned review using the same checkpoint as the proposer, durable
enforcement without sink fencing, and the complete protocol. The first two are
generic controls, not Progent. Each receives the same task permission set and
observable policy state. The learned reviewer must use the named proposer
checkpoint, rather than a stronger hidden service. Per-case hashes bind the
permission and visible-state inputs across paired controller records.

The CPU release measures four scripted controller controls; learned review is
explicitly unrun. In the shared-budget control, every actor has an individually
approved eight-unit request while the shared capacity is eight. Stateless
per-call checks permit each request in isolation; cumulative enforcement admits
only one. This exposes composition failure without suggesting that a stateless
baseline lacks a per-call amount check. The controller workload is sequential;
separate synchronized races measure durable concurrency at 1, 2, 4, 8, 16, and
32 actors. Host results are descriptive and are not pooled into a hardware-free
latency estimate.

The measurement schema includes unauthorized effects, task success,
safe-and-useful completion, false blocks, stranded reservations, fence
acknowledgment delay, and latency. A useful transfer must match recipient and
amount; a useful release must match actual bytes and destination; a useful edit
must match expected contents. The supported synthetic document is scored by
exact expected contents. Arbitrary source-code tasks would need separately
sandboxed behavioral tests; a changed database row alone is insufficient.

False blocks require a reference-permitted task that the controller prevents.
Rejecting the second competing eight-unit request after the budget is exhausted
is not a false block. A cancelled old-epoch operation after an acknowledged fence
is not a newly unauthorized effect. A stranded reservation is an availability
cost, even if it conservatively protects safety. These distinctions are necessary
to avoid rewarding a controller for accepting too much or claiming every
rejection as evidence of successful protection.

\subsection{Public benchmark and upstream reference}
The declared AgentDojo subset contains banking tasks 0, 2, 3, 4, 5, 6, and 11,
and workspace tasks 29, 31, 33, 34, and 35. The upstream repository and benchmark
version are pinned. Native utility and attack graders are retained. Scripted
ground-truth calls supply positive integration controls, and unchanged-state
inputs supply negative utility controls. These establish that the native graders
and selected task environments run; they do not establish model robustness.

Pinned upstream Progent is an additional reference implementation. The current
integration controls provide exact task-call policies to its unmodified policy
checker and execute the native task calls. The subset solver is exercised on
reflexive policy inclusion. This tests tool-argument enforcement and integration
with the selected environments; it does not evaluate Progent's model-generated
policies or adaptive policy updates. Unsupported durable-budget and fencing
properties are listed separately. Import failures, unsupported integrations, and
grader errors would be retained as such, not converted into measured security
failures \cite{agentdojo,progent-code}.

Public benchmark outcomes remain descriptive even after a model run. They are
not fresh procedural holdouts, and the task selection is public. They therefore
do not supply the independent-case inference used for the primary controller
comparisons. Using the native graders prevents an apparently compatible wrapper
from quietly changing what success means.

\subsection{Predeclared model budget and holdouts}
The model manifest names official Qwen3-4B and Qwen3-8B Q4\_K\_M artifacts by
repository, revision, and filename, and pins llama.cpp by commit. Thinking is
disabled. Generation settings, bounded context, maximum generated tokens per
turn, and tool-turn limits are fixed. The two parameter counts identify
checkpoints; they are not a measured ordering of adversarial capability. No
frontier-model extrapolation is authorized by this design.

Total model work has a four-GPU-hour ceiling. Development-case timing may consume
at most 20 minutes. Before inspecting holdout outcomes, the largest sample among
32, 64, and 128 campaigns per checkpoint must fit the remaining time with 20\%
timing headroom. The sample-selection function rejects an infeasible minimum.
One GPU may be used at a time, and active host workloads are never preempted.
Hostname and architecture are verified before any remote allocation. A host
alias resolving to a different machine is not sufficient authorization to run
the job there.

The freeze record binds source, policy, prompt, and analysis inputs before fresh
procedural holdouts are generated. The helper samples fresh randomness only
after checking that record. Repeated variants share a base-case identifier so
their outcomes can be clustered. A changed defense or analysis file invalidates
the freeze. The managed runner checks server and GGUF hashes, launches one local server at a time, enforces phase and total deadlines, and drives the five-controller model-to-tool bridge. CPU scripted controls test that bridge and its learned-review path. Actual GPU execution remains unvalidated. Public-reference work reserves three measured campaign equivalents per checkpoint before sample selection. Its upstream Progent adapter uses generated initial policies and local inference transport; adaptive updates are explicitly disabled and reported. No holdout was generated or inspected
for this release.

\subsection{Inference and noninferiority}
Procedural results must be reported separately for each checkpoint. Controller
comparisons pair by base case and campaign variant; repeated variants are
averaged within a base case. The predeclared analysis uses distribution-free
Hoeffding intervals for independent base-case means, retaining all variants.
This is conservative and avoids treating repeated versions of one case as
independent campaigns. It also avoids a degenerate bootstrap interval implying
certainty when every observed outcome is identical.

For a bounded metric $X\in[0,1]$ over $n$ independent base-case means,
the two-sided radius is $\sqrt{\log(2/\alpha)/(2n)}$. Paired differences
lie in $[-1,1]$ and have twice that radius. The eight primary contrasts are
complete minus each of four baselines at each of two checkpoints. Bonferroni
adjustment uses $\alpha=0.05/8$ for simultaneous primary bounds. This rule is
fixed before a model outcome is available.

The usefulness noninferiority margin is five percentage points. The claim is
made only if the adjusted lower bound for complete-minus-baseline task success
exceeds $-0.05$. A positive point estimate, overlapping marginal intervals, or
an inconclusive bound is insufficient. Unauthorized-effect comparisons and
safe-and-useful completion are reported alongside usefulness, so a controller
cannot appear successful by blocking every task. Missing controllers remain
missing. The code rejects scripted data masquerading as checkpoint outcomes,
duplicate campaign/controller records, mismatched policy-state hashes, and
unpaired variants.
